\raggedbottom % Allow flexible page heights to reduce underfull vbox warnings

\chapter*{Résumé en Français} % Main chapter title
\addcontentsline{toc}{chapter}{Résumé en Français} % Add to table of contents

La cellule constitue l'unité fondamentale du vivant, et la compréhension de ses mécanismes de régulation représente un enjeu majeur de la biologie moderne. Les réseaux de régulation génique (GRN) modélisent les interactions moléculaires qui gouvernent l'expression des gènes au sein de chaque cellule, influençant sa différenciation, sa réponse aux stimuli et son homéostasie. L'inférence précise de ces réseaux à partir de données transcriptomiques unicellulaires demeure toutefois un défi considérable en biologie computationnelle, en raison de l'explosion combinatoire des interactions possibles, du bruit inhérent aux mesures, et de la rareté des ground truths validées expérimentalement. Dans ce contexte, les modèles fondamentaux basés sur les transformers offrent une approche prometteuse pour apprendre des représentations cellulaires riches à partir de grandes quantités de données.

Cette thèse présente des avancées fondamentales dans l'application de l'apprentissage profond basé sur les transformers aux données de séquençage ARN unicellulaire (scRNA-seq). Elle s'articule autour de trois publications principales, chacune abordant un aspect complémentaire : l'inférence de réseaux de régulation génique, l'apprentissage multi-échelle entre modèles biologiques, et la validation systématique des choix architecturaux pour les modèles fondamentaux unicellulaires ainsi qu'un modèle qui bat l'état de l'art sur de nombreuses tâches de biologie cellulaire.

\section*{Chapitre 1 : scPRINT — Pré-entraînement sur 50 millions de cellules pour l'inférence de réseaux géniques}

Le premier chapitre présente scPRINT (single-cell PRe-trained Inference of Networks with Transformers), un modèle cellulaire de grande taille conçu pour l'inférence de réseaux géniques spécifiques à chaque cellule à l'échelle du génome. scPRINT a été entraîné sur plus de 50 millions de cellules provenant de la base de données CellxGene, représentant environ 80 milliards de tokens couvrant plusieurs espèces, maladies et origines ethniques. L'objectif initial de cette thèse était d'utiliser des réseaux de neurones graphiques (GNN) pour améliorer l'inférence de GRN. Cependant, nous avons rapidement constaté que les GNNs ne constituaient pas l'approche optimale : l'absence de structures graphiques connues a priori et leurs propriétés de mise à l'échelle limitées nous ont conduits à adopter les transformers, qui peuvent être vus comme des GNN opérant sur des graphes entièrement connectés.

Le modèle introduit plusieurs innovations architecturales majeures. Premièrement, un encodage génique basé sur les protéines utilisant les embeddings du modèle de langage protéique ESM2, ce qui réduit le nombre de paramètres tout en permettant la généralisation inter-espèces. Deuxièmement, une tokenisation apprise de l'expression génique via des perceptrons multicouches (MLP), qui surpasse les méthodes de discrétisation manuelles utilisées par les modèles existants. Troisièmement, un encodage positionnel génomique qui capture les patterns de co-régulation entre gènes chromosomiquement proches. Le pré-entraînement repose sur trois tâches complémentaires : une tâche de débruitage par suréchantillonnage de transcrits, une tâche d'apprentissage par goulot d'étranglement pour la compression et la reconstruction d'embeddings cellulaires, et une tâche de prédiction de labels avec classification hiérarchique pour obtenir des embeddings cellulaires disentangled représentant différentes facettes phénotypiques.

Une contribution critique est notre méthode d'extraction de réseaux géniques spécifiques à chaque cellule à partir des matrices d'attention du transformer, inspirée d'approches similaires dans ESM2 pour la prédiction de contacts protéiques. Nous avons rendu cette approche extensible pour calculer des réseaux à l'échelle du génome pour des milliers de cellules sur du matériel grand public. Nous avons également introduit un mécanisme de sélection des têtes d'attention, où un sous-ensemble de têtes est choisi en fonction de leur corrélation avec des réseaux de vérité terrain connus, améliorant significativement la qualité des réseaux inférés.

Pour remédier au manque d'évaluation standardisée dans le domaine, nous avons créé BenGRN et GRnnData, des suites de benchmarking pour l'inférence de GRN utilisant plusieurs types de validations : des réseaux basés sur la littérature (Omnipath), intersections ChIP-seq/perturb-seq spécifiques aux types cellulaires, et données perturb-seq à l'échelle du génome. Nos résultats démontrent que scPRINT surpasse toutes les autres méthodes — incluant scGPT, Geneformer v2, DeepSEM et GENIE3 — sur la majorité des benchmarks. Sur le benchmark Omnipath couvrant 26 types cellulaires, scPRINT a retrouvé 67\% plus de connexions que GENIE3 et a montré un enrichissement supérieur pour les facteurs de transcription et leurs cibles validées par ENCODE (20\% des facteurs de transcription avec enrichissement significatif, contre 0\% pour scGPT).

Au-delà de l'inférence de réseaux, scPRINT démontre des performances compétitives même sans ré-entraînement sur des tâches orthogonales : débruitage d'expression (égalant les méthodes de l'état de l'art comme MAGIC et KNNsmoothing2, et les surpassant sur les types cellulaires rares) ; classification de types cellulaires (62\% de précision sans ré-entraînement sur plus de 200 types cellulaires, surpassant les méthodes basées sur les gènes marqueurs comme CellTypist) ; et correction d'effets de lots (obtenant des scores compétitifs sans utiliser de labels de lot). Ces capacités émergent naturellement de l'apprentissage de représentations cellulaires de qualité. L'application à un atlas de 83 000 cellules de tissu prostatique normal et pré-cancéreux (HBP) a permis d'identifier des marqueurs précoces du microenvironnement tumoral dans des lymphocytes B à mémoire, ainsi que des hubs de gènes différentiellement présents dans les fibroblastes associés au HBP. Une étude approfondie de ces réseaux révèle des voies interconnectées liant l'échange ionique, le remodelage de la matrice extracellulaire, le stress oxydatif et l'inflammation chronique : des signatures d'états prémalins. Ce chapitre a conduit à une publication dans Nature Communications.

\section*{Chapitre 2 : Xpressor — Vers des modèles fondamentaux apprenant à travers les échelles biologiques}

Le deuxième chapitre présente Xpressor, un cadre architectural et d'entraînement permettant l'apprentissage inter-échelle entre modèles fondamentaux biologiques. Ce travail adresse une limitation importante : les modèles fondamentaux existant à différentes échelles biologiques — molécules, séquences, cellules, tissus — fonctionnent de manière isolée, incapables d'exploiter le flux d'information entre les échelles. Nous proposons que le vocabulaire de chaque échelle peut être vu comme une représentation compressée de l'échelle inférieure : les acides aminés sont composés d'atomes, les gènes sont définis par des acides aminés, les cellules par l'ensemble de leurs gènes, et les tissus par l'organisation de leurs cellules.

L'architecture Xpressor introduit un mécanisme de compression basé sur l'attention croisée qui transforme les représentations de haute dimension en vecteurs de dimension réduite. Des blocs transformer additionnels effectuent une attention croisée entre les embeddings de sortie d'un modèle fondamental et un ensemble de tokens latents appris, créant un goulot d'étranglement qui compresse $m$ tokens de dimension $d_c$ en $n$ tokens cellulaires de dimension $d_t$, où $n \ll m$ et $d_t < d_c$. Ce même transformer peut ensuite décompresser ces représentations par attention croisée avec des tokens de position (ou d'identité). De plus, l'espace latent du goulot est régularisé par des pertes contrastives et des classificateurs spécifiques à chaque token, garantissant que chaque token capture des informations biologiques distinctes.

Notre seconde contribution est une méthode d'ajustement de modèles d'échelle inférieure lors de l'entraînement de modèles d'échelle supérieure. Nous démontrons ceci en utilisant ESM2 (modèle de langage protéique) comme modèle d'échelle inférieure et scPRINT comme modèle d'échelle supérieure. Plutôt que d'utiliser des embeddings ESM2 figés comme tokens géniques, nous ajoutons un adaptateur MLP entraînable qui transforme chaque embedding protéique durant le pré-entraînement de scPRINT.

Les résultats empiriques montrent des améliorations substantielles : la précision de prédiction des types cellulaires augmente de 0,60 à 0,72 (+20\%) avec l'architecture Xpressor par rapport au pooling standard, et la qualité des embeddings s'améliore de 0,48 à 0,52 (+8\%). L'ajustement multi-échelle améliore également la prédiction des types cellulaires de 0,60 à 0,70 (+17\%) et l'inférence de réseaux géniques sur le benchmark Omnipath de 2,0 à 2,4 EPR (+20\%). Ce chapitre a conduit à un poster à ICML 2025 et est actuellement en révision pour publication dans Bioinformatics Advances.

\section*{Chapitre 3 : scPRINT-2 — Vers la nouvelle génération de modèles fondamentaux cellulaires et de benchmarks}

Le troisième chapitre présente scPRINT-2, un modèle fondamental unicellulaire de nouvelle génération dont les décisions de conception ont été systématiquement validées via un cadre d'évaluation additive sans précédent. Ce travail comble une lacune critique : alors que de nombreux modèles fondamentaux unicellulaires ont été proposés, l'importance relative de leurs choix architecturaux, de leurs stratégies d'entraînement et de leurs modalités de données n'avait auparavant jamais été rigoureusement évaluée de manière isolée.

Nous avons conçu un benchmark complet évaluant 42 configurations différentes de composants, certaines variantes étant entraînées avec différentes initialisations pour générer des bornes d'erreurs statistiques. Cette évaluation a révélé plusieurs résultats clés : le débruitage est supérieur au masquage comme tâche de pré-entraînement ; l'expression non normalisée surpasse l'entrée normalisée ; les tokens géniques basés sur ESM surpassent les embeddings appris de zéro ; l'encodage de localisation génomique améliore la convergence du modèle ; la perte MSE surpasse la ZINB en moyenne, mais une perte hybride ZINB+MSE offre le meilleur équilibre ; et la taille du modèle corrèle avec l'amélioration de l'inférence de réseaux et de la prédiction de types cellulaires.

Pour l'entraînement, nous avons assemblé le plus grand corpus unicellulaire à ce jour : plus de 350 millions de cellules provenant de 16 organismes eucaryotes couvrant plus d'un milliard d'années d'évolution, intégrant des données de CellxGene, du jeu de données Tahoe-100M et de la base scBasecount (20 000 jeux de données GEO retraités), totalisant 25 To de données uniques avec environ 400 000 gènes distincts et 4 764 labels cellulaires différents répartis en 140 000 groupes. Nous avons démontré que la diversité des états cellulaires et la qualité des données importent davantage que le simple nombre de cellules : réduire à 200 jeux de données humains n'a causé qu'une perte minimale de performance, tandis que l'utilisation exclusive de jeux de données à faible diversité a provoqué un effondrement des performances. Nous avons introduit des stratégies d'échantillonnage pondéré par cluster et par qualité pour traiter efficacement ces données hétérogènes.

scPRINT-2 incorpore 12 innovations distinctes validées par notre benchmark, incluant l'architecture XPressor, un encodeur d'expression basé sur GNN exploitant l'information de voisinage, l'attention criss-cross (un mécanisme d'attention sous-quadratique inspiré des Recurrent Interface Networks), une compression basée sur VAE avec pertes de dissimilarité entre tokens cellulaires, et une perte de classification hiérarchique actualisée pénalisant les prédictions en fonction de la distance ontologique plutôt que de la simple exactitude binaire.

Sur le benchmark Open Problems, scPRINT-2 a atteint 75\% de précision en classification zero-shot de types cellulaires, surpassant scPRINT-1 (47\%) et tous les autres modèles fondamentaux zero-shot (40–60\%). Avec notre méthode XPEFT, scPRINT-2 a surpassé toutes les méthodes supervisées et non supervisées existantes sur cette plateforme. Pour le débruitage d'expression, scPRINT-2 est devenu état de l'art, surpassant MAGIC dans tous les contextes testés. Pour l'intégration de lots, les performances zero-shot de scPRINT-2 ont dépassé toutes les autres méthodes.

La capacité de généralisation du modèle au-delà de sa distribution d'entraînement constitue un résultat particulièrement marquant. Sur des données de transcriptomique spatiale Xenium (une modalité absente de l'entraînement), scPRINT-2 a réussi à débruiter l'expression, à imputer 5 000 gènes non observés avec des scores de corrélation comparables aux gènes débruités, et à produire des prédictions biologiquement significatives. Sur des tissus pulmonaires de chat et de tigre (organismes non vus durant l'entraînement), scPRINT-2 a atteint une précision suffisante sur 500 types cellulaires pour parfois corriger les annotations d'experts, avec une précision atteignant 95\% après ajustement fin via XPEFT. L'architecture XPressor a également permis la génération contrefactuelle : en remplaçant les embeddings cellulaires spécifiques à l'organisme de cellules de souris par des embeddings humains, nous avons généré des profils d'expression «~humanisés~» avec un enrichissement de 58\% dans les gènes différentiellement exprimés correctement prédits. Nous avons en outre montré que des embeddings géniques biologiquement significatifs — enrichis pour des voies biologiques connues plutôt que pour de simples valeurs d'expression — n'émergent qu'avec un support architectural explicite pour la compression et la reconstruction. Ce chapitre a conduit à une prépublication actuellement en révision à Nature Methods.

\section*{Autres contributions et perspectives}

Au-delà de ces trois chapitres, cette thèse a donné lieu à de nombreuses contributions additionnelles. Nous avons publié sept paquets Python en open source : scPRINT et scPRINT-2 (modèles et scripts d'entraînement), scDataloader (chargement efficace de milliers de jeux de données unicellulaires avec échantillonnage pondéré sur des milliards d'éléments), BenGRN (benchmark pour l'inférence de GRN), GRNNdata (manipulation conjointe de réseaux et de données unicellulaires), Xpressor (reproduction des expériences du chapitre 2), Simpler Flash (mécanismes d'attention efficaces incluant notre attention flash criss-cross), et Hierarchical Classifier (classification hiérarchique pour données unicellulaires). Des efforts d'accessibilité importants ont été réalisés, incluant des tutoriels sur Google Colab, le déploiement sur le hub de modèles Chan–Zuckerberg et la plateforme Superbio.ai, ainsi qu'une implémentation Docker pour le benchmarking sur la plateforme Open Problems. Des efforts de vulgarisation ont été menés via des articles pour l'Institut Pasteur et le CNRS, des vidéos YouTube, plus de 25 présentations invitées, et cinq posters dans des conférences internationales.

En perspective, les principaux défis identifiés concernent la diversité et la qualité des données (diversité génétique, technologies de séquençage plus profondes), l'intégration de modalités supplémentaires (ATAC-seq, protéomique, transcriptomique spatiale) et de données de perturbation à l'échelle des modèles fondamentaux. À terme, l'état de l'art se dirigera vers la construction d'un modèle cellulaire virtuel piloté par l'IA, combinant pré-entraînement multi-échelle, apprentissage par renforcement avec retour actif, et raisonnement via des modèles de langage, dans une boucle « laboratoire dans la boucle » reliant expériences in silico et in vitro.
